\chapter{Introduction}
\label{chap:introduction}

\section{The Causal-Estimator Evaluation Problem}
\label{sec:introduction:causal-estimator-evaluation-problem}

Causal inference concerns outcomes under alternative actions or exposure conditions, whereas each individual observational record reveals only the factual outcome under the condition that occurred.  The counterfactual outcome needed to score an individual effect estimate is therefore unobserved.  Unlike ordinary prediction, for which a held-out factual outcome can provide a direct target, observational data do not supply a direct causal answer key for evaluating an estimator.  They also leave the true exposure-assignment process and the extent of unmeasured confounding unknown.

Evaluation consequently spans settings with different balances of control and observational fidelity.  Fully synthetic studies specify covariates, exposure assignment, potential outcomes, confounding, support, and effect heterogeneity.  Semi-synthetic studies retain selected empirical components while simulating assignments, outcomes, or both.  Experimentally anchored and empirically anchored designs preserve additional observed structure while obtaining an answer through a controlled experiment, sampling rule, comparison construction, or generated mechanism.  These settings permit causal estimation error to be measured against a specified target, but the evaluator also determines important aspects of the difficulty.  Greater control makes a causal answer measurable, while greater observational fidelity makes the causal target less directly knowable.  These are complementary evaluation properties rather than a simple hierarchy.

CliniCause addresses the observational end of this spectrum.  It retains source-recorded measurements, measurement availability and missingness, care-process traces, covariates, and outcomes while making representation-defined cohorts, proxy prevalence, temporal cutoffs, aggregation, empirical support, DAGs, and adjustment assumptions explicit.  Because no counterfactual truth is revealed, the resources characterize estimator behavior and fragility through non-equivalent evidence rather than scoring error against a known effect.  Chapter~\ref{chap:background-related-work} situates this role relative to controlled causal benchmarks.

\section{Observational Fidelity in Irregular ICU Records}
\label{sec:introduction:motivation-irregular-icu-time-series-analysis}

Intensive-care records combine laboratories, vital signs, treatments, administrative fields, and outcomes observed at nonuniform times and for different purposes.  Variable-specific sampling rates, pervasive missingness, and measurement intensity shaped by clinical concern or workflow make these records high dimensional and heterogeneous.  Irregular medical time-series methods must therefore represent timing, absence, and changing observation availability while preserving a defensible link to the analytical question \cite{sun_2026_review_irregular_medical_timeseries}.

Regular resampling, imputation, and summary aggregation choose a temporal resolution, which observations count, and how absence is represented.  Sparse-event or missingness-aware models preserve different aspects of the observation process, but exploiting masks, gaps, or measurement intensity predictively does not remove their causal implications.  Missingness and irregularity can affect the representation, cohort, rule-derived labels, and adjustment; this thesis therefore treats them as motivations for explicit design and limitations, not as local causal findings.

Two complementary needs follow.  First, an analysis must accommodate irregular longitudinal measurements rather than force all observations into an unqualified regular grid.  Architectures for sparse clinical series and missingness-aware recurrent models illustrate different ways to use values, time intervals, and observation patterns \cite{tipirneni2022strats,che2018grud}.  Second, the analysis needs interpretable variables that can connect a rich event stream to later questions about prognosis and adjusted mortality contrasts.  In this thesis, that intermediate layer consists of \emph{clinically inspired proxy states}: rule-derived analytical constructs that summarize evidence for selected clinical patterns from available measurements.

This wording is deliberate.  A rule-derived proxy label is not a chart-adjudicated diagnosis, ground truth, or validated phenotype: a positive tag records that project-specific rules fired on available data, while a negative tag does not establish absence of a condition.  The transparent proxy layer supplies multi-label prediction targets and documented observational exposures while retaining a visible link to their measurements and rules.  Electronic-phenotyping research provides context for this design choice but does not validate the local constructs \cite{banda_2018_electronic_phenotyping}.

Prediction and causal analysis then answer different questions.  Predicting a rule-derived proxy label asks whether an irregular-time-series model can reproduce that label from the available inputs.  Estimating an adjusted mortality contrast for a proxy-state exposure asks a separate observational question about modeled differences in outcome under explicit assumptions.  Strong performance on the first task does not establish the clinical validity of the label, causal validity of the second analysis, or actionability of the exposure.  The latter task requires a stated causal model, project-specified directed acyclic graphs (DAGs), exposure-specific adjustment logic, attention to support and overlap, sensitivity and permutation diagnostics, and traceable evidence.  Even with these elements, adjusted estimates remain conditional on the adequacy of the proxy measurement, graph, temporal ordering, observed covariates, and model specification \cite{pearl_1995_causal_diagrams,hernan_robins_2016_target_trial,hernan_taubman_2008_well_defined_interventions}.

The framework is instantiated separately on PhysioNet 2012 and MIMIC-III \cite{silva2012physionet,johnson2016mimiciii}.  Their eras, scales, variables, data contracts, measurement processes, proxy definitions, and project graphs differ, precluding automatic construct equivalence or numerical pooling.  Explicit handoffs from source-specific contracts through proxy construction and prediction to adjusted analysis expose whether a difference entered through the record, representation, export, or estimator.

\section{Representation-Centered Observational Testbeds}
\label{sec:introduction:representation-centered-observational-testbeds}

The central move in CliniCause is to shift controlled researcher intervention away from simulating the patient-level data-generating process and toward specifying an inspectable \emph{representation layer}.  That layer comprises proxy constructs, operational rules and temporal cutoffs, analytical cohorts, aggregation, project-specified DAGs, exposure-specific adjustment assumptions, and provenance.  Patient records contribute source-recorded measurements, availability and missingness patterns, care-process traces, covariates, and in-hospital mortality; empirical support and proxy prevalence emerge only after the representation and cohort are defined.  CliniCause instantiates proxy labels deterministically and does not simulate records, exposure-assignment mechanisms, or outcomes.  Source recording does not imply error-free or causally sufficient data.

The design and instantiation roles are deliberately separated.  The LLM received schema-level information---available column names and the mortality outcome---together with instructions to use relevant literature.  It did not receive patient records, empirical distributions, exposure prevalences, estimated effects, or execution access.  Its output consisted of literature-grounded design proposals for proxy constructs, operational rules, missingness considerations, and dataset-specific graphs.  Project-selected rules and graphs were subsequently encoded as deterministic source.  The LLM was therefore a design assistant, not a clinical or causal expert, an executed patient-level component, or a source of ground truth.

This representation is instantiated independently on MIMIC-III and PhysioNet 2012.  The resulting resources have a common analytical role but are not pooled replicas of one object.  Similarly named proxy constructs can differ in measurement availability, operational definition, prevalence, missingness, support, graph structure, adjustment, and population.  Each resource is a \emph{testbed} when used to study causal-analysis behavior; the term does not imply a benchmark with known causal truth.

Figure~\ref{fig:clinicause-testbed-pipeline} summarizes the separation among design, deterministic instantiation, and evaluation.

\begin{figure}[htbp]
  \centering
  \includegraphics[width=\textwidth]{clinicause_testbed_pipeline.png}
  \caption[CliniCause design, instantiation, and evaluation stages.]{The Design, Instantiate, and Evaluate stages of CliniCause.  Schema and literature information support LLM-assisted representation proposals, which project-selected deterministic code instantiates at patient level on source-recorded ICU records and outcomes.  The resulting dataset-specific resources are characterized through converging evidence without a known causal answer key; the LLM receives schema information rather than patient records.}
  \label{fig:clinicause-testbed-pipeline}
\end{figure}

\section{Thesis Gap and Objective}
\label{sec:introduction:thesis-gap-objective}

Prior work supplies the components rather than the complete evaluation object: electronic phenotyping and data programming motivate inspectable analytical labels \cite{banda_2018_electronic_phenotyping,ratner_et_al_2016_data_programming}; medical-LLM and graph-prior studies motivate bounded design proposals \cite{singhal_et_al_2023_llm_clinical_knowledge,darvariu_et_al_2024_llm_causal_graph_priors}; and causal-inference work supplies DAG, DML, forest, and ICU-design principles \cite{chernozhukov2018dml,wager2018causalforest,smit_2023_causal_inference_icu_scoping_review}.  None validates the local proxy rules, graphs, or estimates.

The thesis objective is to construct and characterize clinically grounded observational causal-analysis testbeds from irregular, source-recorded ICU data.  Its evidence-tracked representation framework makes proxy definitions, data contracts, causal assumptions, analytical interfaces, diagnostics, provenance, and limitations explicit.  Evidence tracking distinguishes design proposals, selected source behavior, checked numerical artifacts, and unresolved validation or lineage requirements.

At design time, a structured LLM-assisted protocol elicited proxy ontologies, rule families, missingness considerations, and dataset-specific DAG proposals; project-selected designs were encoded as deterministic source.  Execution then applies source-specific preprocessing, rule-label construction, multi-label prediction, normalization and fixed aggregation, graph-guided adjustment, matching and effect estimation, and robustness/provenance evaluation.  The LLM is not invoked on patient records or during estimation.

These interfaces make the testbeds inspectable without upgrading their evidence.  Prediction metrics concern rule-derived targets; matching is a descriptive comparison; mean model-estimated CATE is conditional on the implemented representation and assumptions; and support, permutation, and sensitivity diagnostics address particular failure modes.  The framework provides neither clinical validation nor unconditional causal effects, recommendations, deployment readiness, or full historical rerun reproducibility.

\section{Research Questions and Contributions}
\label{sec:introduction:research-questions-contributions}

The main research question is:

\begin{quote}
How can clinically grounded observational causal-analysis testbeds be constructed from source-observed irregular ICU records---without simulating patient records, treatment or exposure-assignment mechanisms, or outcomes---and what converging evidence can characterize their usefulness and limitations for evaluating causal estimators?
\end{quote}

The words ``observational'' and ``characterize'' are central.  The thesis investigates how a linked workflow can construct and qualify retrospective evaluation resources; it does not claim to score estimates against causal truth, prove that proxy-state exposures cause mortality, or show that manipulating a proxy would improve outcomes.  The main question is addressed through five secondary research questions:

\begin{description}
  \item[SRQ-1:] How can clinically and literature-grounded representation designs be instantiated deterministically and separately in MIMIC-III and PhysioNet while retaining source-specific measurement and missingness processes?
  \item[SRQ-2:] Which data contracts, normalization, aggregation, DAG, adjustment, and provenance mechanisms make the constructed resources explicit, auditable, and reusable?
  \item[SRQ-3:] How recoverable are the deterministic proxy labels from irregular time series, and what mortality-relevant predictive information is retained by the constructed proxy representations?
  \item[SRQ-4:] What converging evidence---including estimator sign, rank, and magnitude agreement, clinical-literature comparison, matching and support, permutation disruption, and omitted-variable sensitivity---characterizes each resource?
  \item[SRQ-5:] Which limitations prevent the resources from being interpreted as clinical validation, known causal truth, or proof of estimator correctness?
\end{description}

The primary contribution is a representation-centered framework and two concrete, dataset-specific observational causal-analysis testbeds constructed from source-recorded MIMIC-III and PhysioNet 2012 records.  The framework keeps proxy definitions, operational rules, cohort choices, DAGs, adjustment assumptions, and provenance explicit and replaceable; converging, non-equivalent evidence characterizes the resources because no causal answer key is available.

Enabling methods span schema-bounded LLM-assisted elicitation, deterministic proxy and DAG source, dataset-specific contracts, proxy prediction and aggregation, adjustment, matching, estimation, diagnostics, and provenance.  CausalForestDML is primary, LinearDML secondary, and CausalPFN exploratory because it lacks the DML estimators' diagnostic support.

Empirically, four prediction models characterize proxy recoverability; original-cohort analyses retain all 19 dataset--exposure combinations; and outcome-downsampling and cross-resource comparisons remain separate, without pooling.  Forest and Linear agree in direction for 19 of 19 combinations, while all three estimators agree for 18 of 19 and retain PhysioNet shock as the exception.  Matching/support, mortality prediction, clinical comparison, disruption, omitted-confounding sensitivity, population perturbation, and provenance supply additional qualified evidence.  Checked tables, manifests, hashes, and decision records support numerical traceability but not a demonstrated clean-checkout historical rerun.

These contributions are cumulative, not interchangeable: workflow integration makes handoffs inspectable, empirical comparisons characterize the resources, and evidence records separate numerical traceability from reproducibility.  Agreement does not establish estimator accuracy, graph correctness, identification, or clinical validity, and no layer upgrades a proxy label to a diagnosis or a graph-guided contrast to a well-defined intervention effect.  STraTS led the archived MIMIC metrics and GRU-D the PhysioNet metrics, a resource-dependent result whose detailed qualifications appear in Chapter~\ref{chap:results}.

\section{Thesis Organization}
\label{sec:introduction:thesis-organization}

The remainder of the thesis follows construction and characterization of the two testbeds from context to interpretation.  Background and study design establish the evaluation spectrum, resource contracts, questions, and assumptions.  Chapters~\ref{chap:data-preprocessing}--\ref{chap:experimental-design} construct and characterize the resources through source-specific contracts, proxy instantiation and prediction, DAG-guided adjustment and estimation, diagnostics, and provenance.  The final chapters report the checked evidence, interpret its limitations, and state the conclusions and future-work priorities.
